%==============================================================================
% SLHA v2 -- arXiv preprint
% Author: Tarek Zekriti (ZEKRITI TAREK)
% Self-contained: compiles with pdflatex (no external .bib, no custom .sty).
%   pdflatex slhav2.tex ; pdflatex slhav2.tex
%==============================================================================
\documentclass[11pt]{article}

\usepackage[utf8]{inputenc}
\usepackage[T1]{fontenc}
\usepackage{lmodern}
\usepackage{microtype}
\usepackage{amsmath,amssymb,amsthm}
\usepackage{booktabs}
\usepackage{array}
\usepackage{multirow}
\usepackage{graphicx}
\usepackage{xcolor}
\usepackage{listings}
\usepackage{enumitem}
\usepackage[margin=1in]{geometry}
\usepackage{caption}
\usepackage[hidelinks]{hyperref}

\hypersetup{
  pdftitle={SLHA v2: Cache-Aware Hybrid Attention with Low-Rank Latents and 1-Bit Sign-LSH Residuals},
  pdfauthor={Tarek Zekriti}
}

% ---- theorem environments ----------------------------------------------------
\theoremstyle{plain}
\newtheorem{lemma}{Lemma}
\newtheorem{proposition}{Proposition}
\theoremstyle{definition}
\newtheorem{definition}{Definition}

% ---- listings (Rust) ---------------------------------------------------------
\definecolor{kw}{rgb}{0.10,0.10,0.55}
\definecolor{cmt}{rgb}{0.30,0.45,0.30}
\definecolor{str}{rgb}{0.55,0.20,0.10}
\lstdefinelanguage{Rust}{
  keywords={fn,let,mut,pub,struct,impl,for,if,else,return,in,const,usize,u8,u16,u32,u64,f32,i32,self,Self,as,bool,enum,match,unsafe},
  sensitive=true,
  morecomment=[l]{//},
  morestring=[b]",
}
\lstset{
  language=Rust,
  basicstyle=\ttfamily\footnotesize,
  keywordstyle=\color{kw}\bfseries,
  commentstyle=\color{cmt}\itshape,
  stringstyle=\color{str},
  showstringspaces=false,
  columns=fullflexible,
  breaklines=true,
  frame=single,
  framesep=4pt,
  rulecolor=\color{black!25},
  aboveskip=8pt,belowskip=6pt,
}

% ---- shorthands --------------------------------------------------------------
\newcommand{\dc}{d_c}
\newcommand{\ds}{d_s}
\newcommand{\dk}{d_k}
\newcommand{\sigE}{\sigma_E}
\newcommand{\inner}[2]{\langle #1, #2 \rangle}
\newcommand{\popc}{\mathrm{popcount}}
\newcommand{\sign}{\mathrm{sign}}
\newcommand{\rms}{\mathrm{RMS}}

\title{\textbf{SLHA v2: Cache-Aware Hybrid Attention with Low-Rank
Latents and 1-Bit Sign-LSH Residuals for Memory-Bandwidth-Bound
LLM Inference}}

\author{
  Tarek Zekriti\thanks{Independent Researcher. Correspondence:
  \texttt{zekrititarek@gmail.com}.}\\
  \small Independent Researcher
}

\date{June 2026}

\begin{document}
\maketitle

\begin{abstract}
\noindent
Autoregressive inference of large language models (LLMs) is bounded by
memory bandwidth: the key/value (KV) cache grows linearly with context
length and saturates the interconnect long before the arithmetic units
are busy. We present \textbf{SLHA~v2}, a cache-aware attention mechanism
that decomposes each key into (i)~a shared low-rank latent base
quantized to 4~bits and (ii)~a deterministic 1-bit sign-LSH residual.
The two are combined in a single \emph{fused} score that adds a
continuous dot product to a branchless Hamming-distance correction whose
exactness we prove. The representation is packed into a hardware-aligned
\textbf{128-byte tile with zero padding} that maps onto cache lines, and
is managed by an elastic KV cache (\emph{CCOS Soft-Paging}) whose
three-state \textsc{hot}/\textsc{warm}/\textsc{cold} machine bounds the
memory footprint under a byte budget \emph{without I/O}, by freeing
residuals where they matter least. We provide an open, zero-dependency
Rust reference implementation (\textsc{SciRust}) with
runtime-dispatched scalar / AVX2 / AVX-512 / NEON kernels proven
equivalent by 37 automated tests. On reproducible synthetic benchmarks
(fixed seeds) we measure: a 1-bit core that exactly realizes the
signed-dot identity; an attention-\emph{output} cosine fidelity of
$0.95$--$0.997$ versus full precision, \emph{despite} score-ranking
Spearman plateauing at $0.79$--$0.90$ (the softmax absorbs score error);
$\sim\!2.5\times$ more scored tokens/s than a \texttt{bf16} baseline at
comparable memory bandwidth; SIMD speedups of $11.5\times$ (AVX2) and
$14.1\times$ (AVX-512); and a learned, task-aware low-rank projection
that lifts \textsc{warm} Spearman from $0.16$ to $0.86$ under query/key
distribution shift. A controlled ablation shows that quantization width
is \emph{not} the bottleneck (an INT8 reference does not beat INT4 on the
coarse term)---the low-rank projection is. We position SLHA~v2 as a
validated \emph{mechanism} and a systems substrate, and we are explicit
about what remains unmeasured: real-model perplexity, hardware cache
counters, and joint projection training.
\end{abstract}

\noindent\textbf{Keywords:} efficient inference, KV cache, low-rank
attention, 1-bit quantization, locality-sensitive hashing, cache-aware
data structures, SIMD.

\section{Introduction}
\label{sec:intro}

Conventional attention pipelines treat the KV cache as a contiguous
stream of high-precision tensors (FP16/BF16). During decoding, each new
token re-reads the entire history from main memory into the registers of
the CPU/GPU. The cost of this re-read is set by the cache hierarchy:

\begin{itemize}[leftmargin=1.4em,itemsep=2pt]
  \item \textbf{L1 (data):} $\sim\!1$ cycle, but tiny (32--64\,KiB/core);
  \item \textbf{L2:} $\sim\!4$--$15$ cycles, $0.5$--$2$\,MiB/core;
  \item \textbf{L3 / LLC:} shared, $\sim\!40$--$80$ cycles, tens to
    hundreds of MiB.
\end{itemize}

If the KV layout does not align with these strata, the processor spends
most of its time waiting on transfers (cache misses) rather than
computing. This is the \emph{memory-bandwidth wall}: for long contexts,
arithmetic intensity collapses and FLOPs become free relative to bytes.

SLHA~v2 attacks the wall at the level of the data structure itself. Its
contributions are:

\begin{enumerate}[leftmargin=1.6em,itemsep=2pt]
  \item \textbf{An asymmetric key representation} (\S\ref{sec:method}):
    a shared low-rank latent (4-bit) plus a deterministic 1-bit
    sign-LSH residual, fused into one score that combines a float dot
    product with a Hamming correction. We prove the binary core is an
    exact signed dot product (Lemma~\ref{lem:hamming}).
  \item \textbf{A hardware-aligned 128-byte tile}
    (\S\ref{sec:tile}) with \emph{zero padding}, whose 64-byte latent
    sub-block is exactly one cache line.
  \item \textbf{An elastic KV cache with Soft-Paging}
    (\S\ref{sec:ccos}): a \textsc{hot}/\textsc{warm}/\textsc{cold}
    state machine that bounds the footprint under a byte budget by
    masking residuals ($O(1)$, no I/O, no allocation) before evicting
    the most causally distant tokens.
  \item \textbf{An open, zero-dependency reference implementation}
    (\S\ref{sec:impl}) with runtime-dispatched SIMD kernels proven
    equivalent to a scalar reference.
  \item \textbf{A reproducible, deliberately honest evaluation}
    (\S\ref{sec:eval}) that separates what is \emph{measured} from what
    is \emph{targeted}, including negative results
    (\S\ref{sec:discussion}).
\end{enumerate}

A guiding finding, established by ablation (\S\ref{sec:ablation}), is
that the fidelity ceiling of the coarse term is set by the
\emph{low-rank projection}, not by latent quantization width: doubling
the bits (INT8) does not help, while a \emph{learned} projection lifts
quality dramatically (\S\ref{sec:learned}). This reframes where future
effort should go.

\paragraph{Scope statement.}
This paper reports the validation of a \emph{mechanism} on synthetic
data with reproducible, seeded benchmarks. We do not claim end-to-end
perplexity on a real model; \S\ref{sec:limits} states the boundaries
precisely. We believe an honest account of a partially validated system
is more useful than an overclaimed one.

\section{Background and Related Work}
\label{sec:related}

\paragraph{Low-rank / latent KV compression.}
Multi-head Latent Attention (MLA) in DeepSeek-V2~\cite{deepseekv2}
projects keys and values through a shared low-rank bottleneck to shrink
the KV cache. SLHA~v2 adopts the same low-rank \emph{base} but adds a
quantized 1-bit residual to recover the information the bottleneck
discards, and ties the layout to the cache line.

\paragraph{KV-cache quantization.}
GPTQ~\cite{gptq}, AWQ~\cite{awq}, and KIVI~\cite{kivi} quantize weights
or KV entries to low bit-width. NF4 (from QLoRA~\cite{qlora}) uses a
normal-float codebook. SLHA~v2 uses signed INT4 (with optional grouped
micro-scaling, an NF4 codebook, and TQ3, a port of the
TurboQuant\footnote{\url{https://github.com/CHECKUPAUTO/TurboQuant}} KV
codec) for the latent, but our ablation
(\S\ref{sec:ablation}) shows quantization width is not the limiting
factor here.

\paragraph{1-bit and salient-aware quantization.}
BiLLM~\cite{billm} pushes post-training quantization toward 1~bit by
isolating salient weights in higher precision. SLHA~v2's residual is
1-bit \emph{by construction} (a sign-LSH code), and \S\ref{sec:future}
discusses a BiLLM-style salient extension as future work.

\paragraph{Locality-sensitive hashing.}
Sign-LSH / SimHash~\cite{charikar} estimates the angle between two
vectors from the agreement of random hyperplane signs; the underlying
random projection is justified by the Johnson--Lindenstrauss
lemma~\cite{jl}. SLHA~v2's residual is a sign-LSH code of the
reconstruction error, and the fused score reads it back via a Hamming
distance.

\paragraph{KV eviction and paging.}
H2O~\cite{h2o} and StreamingLLM~\cite{streamingllm} drop or retain KV
entries by heuristics; PagedAttention/vLLM~\cite{vllm} manages KV memory
in pages. SLHA~v2's Soft-Paging adds an intermediate, near-lossless
\textsc{warm} state---freeing only the 1-bit residual---\emph{before}
resorting to eviction, and chooses what to free by residual energy.

\section{The SLHA v2 Mechanism}
\label{sec:method}

\subsection{Asymmetric decomposition}
SLHA~v2 splits the representation of each key into two asymmetric
components: a compact semantic base and a high-frequency binary
correction. Let $x_j \in \mathbb{R}^{d_{\text{model}}}$ be the
activation of context token $j$.

\subsection{Multi-head latent compression}
A shared down-projection maps $x_j$ into a latent bottleneck of
dimension $\dc$:
\begin{equation}
  h_{KV,j} \;=\; W_{\text{down}}\, x_j \;\in\; \mathbb{R}^{\dc}.
\end{equation}
Each head $n$ reconstructs a \emph{coarse key} via a head-specific
up-projection $W_{\text{up},K}^{(n)} \in \mathbb{R}^{\dk \times \dc}$:
\begin{equation}
  K_{\text{coarse},j}^{(n)} \;=\; W_{\text{up},K}^{(n)}\, h_{KV,j}.
\end{equation}

\subsection{Deterministic 1-bit sign-LSH residual}
The reconstruction error
$E_j^{(n)} = K_{\text{real},j}^{(n)} - K_{\text{coarse},j}^{(n)}$ is
captured by a fixed random orthogonal Johnson--Lindenstrauss projection
$Z \in \mathbb{R}^{\ds \times \dk}$, keeping only the sign on $\ds$
dimensions:
\begin{equation}
  B_j^{(n)} \;=\; \sign\!\left( Z\, E_j^{(n)} \right) \;\in\; \{-1,+1\}^{\ds}.
\end{equation}
$Z$ is stable and pseudo-random (seeded at start-up); $B_j^{(n)}$ is
stored as compact bitmaps. With $\ds = 256$ the residual occupies
exactly 32~bytes (four 64-bit words).

\subsection{Fused float--binary score}
The unnormalized attention score between query $Q_i^{(n)}$ and cached
token $j$ combines a continuous inner product and a binary inner product
accelerated by bitwise logic:
\begin{equation}
  s_{ij}^{(n)} \;=\;
  \underbrace{\inner{Q_i^{(n)} W_{\text{up},K}^{(n)}}{h_{KV,j}}}_{\text{coarse (continuous)}}
  \;+\;
  \lambda \cdot
  \underbrace{\Big[\,\ds - 2\,\popc\big(p(Q_i^{(n)}) \oplus B_j^{(n)}\big)\Big]}_{\text{residual (1-bit)}},
  \label{eq:score}
\end{equation}
where $\oplus$ is bitwise XOR, $\popc$ counts set bits, and
$p(Q_i^{(n)}) = \mathrm{pack\_sign}(Z\,Q_i^{(n)})$ packs the query's
sign-LSH code into the same bitmap format. The next lemma shows the
bracketed term is an \emph{exact} signed dot product, not an
approximation of one.

\begin{lemma}[Hamming identity]
\label{lem:hamming}
Let $a,b \in \{0,1\}^{\ds}$ encode signs via $\sigma(0)=+1$,
$\sigma(1)=-1$, and let $s_a, s_b \in \{-1,+1\}^{\ds}$ be the
corresponding sign vectors. Then
\begin{equation}
  \ds - 2\,\popc(a \oplus b) \;=\; \sum_{k=1}^{\ds} s_{a,k}\, s_{b,k}
  \;=\; \inner{s_a}{s_b}.
\end{equation}
\end{lemma}
\begin{proof}
Bit $k$ of $a \oplus b$ is $1$ iff $a_k \neq b_k$, i.e.\ iff
$s_{a,k} s_{b,k} = -1$. Let $D = \popc(a\oplus b)$ be the number of
disagreements. Agreements contribute $+1$ and disagreements $-1$, so
$\inner{s_a}{s_b} = (\ds - D)\cdot(+1) + D\cdot(-1) = \ds - 2D$.
\end{proof}

This identity is verified against a brute-force reference for $2000$
random bitmap pairs (test \texttt{hamming\_identity\_matches\_sign\_dot},
\S\ref{sec:eval}). It is what makes the residual term a single XOR
followed by a population count---two of the cheapest operations on any
modern ISA.

\subsection{Dynamic $\lambda$ calibration}
\label{sec:lambda}
The weight $\lambda$ of the binary correction is not fixed. From the
residual energy $\sigE = \rms(E_j^{(n)})$ the spec sets
\begin{equation}
  \lambda \;=\; \sigE \cdot \sqrt{\tfrac{\pi}{2\,\ds}}.
  \label{eq:lambda}
\end{equation}
The \emph{shape} $\lambda \propto \sigE$ is validated empirically
(\S\ref{sec:calib}): the closed-form optimal multiplier against a
full-precision reference is near-constant across residual energy. The
analytic \emph{constant} $\sqrt{\pi/(2\ds)} \approx 0.078$ at $\ds=256$
under-weights the residual by $\approx 4.2\times$; the calibrated
constant is $C_{\text{emp}} \approx 0.33$. We keep the analytic form as
the conservative default and document the calibrated constant
(\S\ref{sec:calib}).

\section{Hardware-Aware Tile Layout}
\label{sec:tile}

\subsection{The 128-byte tile}
SLHA~v2 stores each (key, residual, metadata) triple in a single
\texttt{\#[repr(C)]} structure whose total size is exactly 128~bytes
\emph{with zero padding} --- two 64-byte cache lines. It is \texttt{align(64)}
by default, so it never straddles a third line; \texttt{build.rs} raises it to
\texttt{align(128)} on a native 128-byte-line host (\S\ref{sec:cachemap}).
Table~\ref{tab:tile} gives the layout.

\begin{table}[t]
\centering
\caption{The \texttt{SciRustSlhaTile} layout ($\dc=128$, $\ds=256$).
Total $= 64+32+32 = 128$ bytes; \texttt{align\_of}~$=64$ by default (two cache
lines), raised to 128 by \texttt{build.rs} on a native 128-byte-line host;
verified by test \texttt{tile\_is\_exactly\_128\_bytes\_zero\_padding} and by
AArch64 cross-compilation.}
\label{tab:tile}
\small
\begin{tabular}{@{}llr l@{}}
\toprule
Field & Type & Bytes & Role \\
\midrule
\texttt{latent\_kv}       & \texttt{[u8; 64]}  & 64 & 128 INT4/NF4 latent dims (2/byte) \\
\texttt{residual\_bitmap} & \texttt{[u64; 4]}  & 32 & sign-LSH residual (256 bits) \\
\texttt{scale}            & \texttt{f32}       & 4  & global INT4 dequant scale \\
\texttt{dynamic\_lambda}  & \texttt{f32}       & 4  & residual weight $\lambda$ (Eq.~\ref{eq:lambda}) \\
\texttt{residual\_sigma}  & \texttt{f32}       & 4  & per-tile $\sigE$ (for $\lambda$) \\
\texttt{token\_id}        & \texttt{u32}       & 4  & token identifier \\
\texttt{position}         & \texttt{u32}       & 4  & causal position \\
\texttt{head\_id}         & \texttt{u16}       & 2  & attention head \\
\texttt{flags}            & \texttt{u16}       & 2  & \textsc{hot}/\textsc{warm}/NF4 \\
\texttt{group\_scales}    & \texttt{[u8; 8]}   & 8  & per-group MX micro-scales \\
\bottomrule
\end{tabular}
\end{table}

The signed INT4 dequantization uses a baked-in zero point,
$\hat v_d = (\text{nibble}_d - 8)\cdot \texttt{scale}$, so the latent
represents negative values without a separate zero-point array. Optional
grouped micro-scaling (MX) stores one \texttt{u8} scale per group of 16
dims in \texttt{group\_scales}, giving effective scale
$\texttt{scale}\cdot \texttt{group\_scales}[g]/255$.

\begin{lstlisting}[caption={Core tile and fused score (abridged from
\texttt{scirust/src/attention/slha\_v2.rs}).},captionpos=b]
// 128 bytes = two 64-byte cache lines (x86-64 and Neoverse both use 64B lines).
// align(128) only on a native 128B-line host (build.rs sets cfg(cache_line_128)).
#[cfg_attr(cache_line_128, repr(C, align(128)))]
#[cfg_attr(not(cache_line_128), repr(C, align(64)))]
pub struct SciRustSlhaTile {
    pub latent_kv: [u8; 64],         // INT4/NF4 latent (128 dims)
    pub residual_bitmap: [u64; 4],   // sign-LSH residual (256 bits)
    pub scale: f32, pub dynamic_lambda: f32, pub residual_sigma: f32,
    pub token_id: u32, pub position: u32,
    pub head_id: u16, pub flags: u16,
    pub group_scales: [u8; 8],
}

// score = <q_coarse, dequant(latent)>  +  lambda * (d_s - 2*popcount(q_sign ^ B))
// Branchless Hamming reduction; one vpopcntq over 256 bits on AVX-512 VPOPCNTDQ,
// else count_ones() (-> POPCNT/CNT). Runtime-selected; bit-equivalent by test.
pub fn hamming_distance(a: &[u64; 4], b: &[u64; 4]) -> u32 {
    #[cfg(target_arch = "x86_64")]
    if is_x86_feature_detected!("avx512vpopcntdq") && is_x86_feature_detected!("avx512vl") {
        return unsafe { hamming_vpopcntdq(a, b) };
    }
    (0..4).map(|w| (a[w] ^ b[w]).count_ones()).sum()
}
\end{lstlisting}

\subsection{Cache mapping}
\label{sec:cachemap}
The 64-byte latent sub-block is exactly one cache line and holds all 128
semantic dimensions; it is this sub-block---not the whole tile---that
feeds the arithmetic unit in a single line load. Queries are pre-loaded
into L1 and the context tiles are swept sequentially, so the access
pattern is a pure streaming read with unit stride.

Both of our target families use a \textbf{64-byte} cache line: x86-64, and
AArch64/Neoverse-V3AE. We initially assumed a Jetson Thor-class part would
use 128-byte lines (the ``128'' in \emph{AGX 128}), but measuring the device
refuted this---all levels are 64~bytes (L1d/L1i/L2; see Table~\ref{tab:thor})
---so the ``128'' denotes the 128~GB \emph{unified CPU/GPU} LPDDR5X memory,
not the cache line.
We therefore default the tile to \texttt{align(64)}, which is correct and
optimal on both families: a 64-aligned 128-byte tile occupies exactly two full
lines and never straddles a third, and \texttt{align(64)} prevents false
sharing without the dead space a larger alignment would imply. A part that
genuinely uses 128-byte lines (e.g.\ Apple Silicon) benefits from
\texttt{align(128)} for a single-line fill; because alignment must be a
compile-time constant, the shipped \texttt{build.rs} detects the \emph{host's}
true L1d line size on a \emph{native} build (host triple $=$ target triple;
Linux \texttt{sysfs} or macOS \texttt{sysctl}) and raises the alignment via
\texttt{cfg(cache\_line\_128)} only there --- never as an \texttt{aarch64}-wide
assumption (\S\ref{sec:future}). For cross-compilation the host line size is
irrelevant, so the 64-byte default holds. We cross-compile to
\texttt{aarch64-unknown-linux-gnu} in CI so the AArch64 layout and NEON kernel
are type-checked.

\section{CCOS Elastic KV-Cache and Soft-Paging}
\label{sec:ccos}

SLHA~v2's decomposition lets the runtime apply a \emph{boundedness}
policy as a fine-grained elasticity of fidelity rather than all-or-nothing
eviction. Each tile occupies one of three states
(Table~\ref{tab:states}).

\begin{table}[t]
\centering
\caption{Soft-Paging states. Footprint is the \emph{logical} byte cost
used by the budget.}
\label{tab:states}
\small
\begin{tabular}{@{}llr@{}}
\toprule
State & Behavior & Footprint \\
\midrule
\textsc{hot}  & full fused score (latent + 1-bit residual) & 128 B \\
\textsc{warm} & residual freed ($\lambda=0$, \texttt{FLAG\_WARM}); score $=$ coarse only & 96 B ($-25\%$) \\
\textsc{cold} & evicted from the working set; slot recycled on next insert & 0 B \\
\bottomrule
\end{tabular}
\end{table}

\subsection{State machine and budget enforcement}
A reference manager, \texttt{ccos::ElasticKvCache}, holds tiles in a
contiguous arena. \texttt{enforce\_budget()} bounds the logical footprint
under a byte budget in two phases:

\begin{enumerate}[leftmargin=1.6em,itemsep=2pt]
  \item \textbf{Page} \textsc{hot}$\to$\textsc{warm} in
    \texttt{PageOutPolicy} order. The default \emph{hybrid} policy pages
    the lowest-$\sigE$ tiles first---their 1-bit residual matters least,
    so \textsc{warm} is near-lossless there (\S\ref{sec:scorefp}).
    Masking is $O(1)$: zero the 32-byte bitmap and set a flag, no
    allocation, no I/O.
  \item \textbf{Evict} live tiles $\to$\textsc{cold} if still over
    budget, \emph{always oldest-first} (causal distance). Dropping a
    token is the harder loss, so it targets the most distant context
    regardless of $\sigE$. The freed slot is recycled via a free list.
\end{enumerate}

This two-key design (page by energy, evict by age) is pinned by tests
(\S\ref{sec:eval}). A production CCOS would snapshot \textsc{cold} tiles
to an append-only event log; we do not simulate that persistence here.

\subsection{Soft-Paging fidelity}
The example \texttt{ccos\_softpaging} streams 8192 tiles (1024\,KiB
all-\textsc{hot}) and measures the effect on the attention
\emph{output} (Table~\ref{tab:softpaging}). Paging \emph{half} the
tiles---the lowest-$\sigE$ ones---to \textsc{warm} leaves the output at
cosine $\approx 0.9995$ versus all-\textsc{hot}: the central benefit of
Soft-Paging is to bound memory by freeing the residual where it weighs
least, with no I/O and no dropped token.

\begin{table}[t]
\centering
\caption{Soft-Paging under a byte budget (8192 tiles, $d_v=64$).}
\label{tab:softpaging}
\small
\begin{tabular}{@{}lllll@{}}
\toprule
Regime & Budget & \textsc{hot}/\textsc{warm}/\textsc{cold} & Footprint & $\cos(\text{out}, \text{all-}\textsc{hot})$ \\
\midrule
A --- paging only      & 896\,KiB (112 B/tile) & 4096 / 4096 / 0    & 896\,KiB (88\%)      & $\mathbf{0.9995}$ \\
B --- forced eviction  & 320\,KiB (40 B/tile)  & 0 / 3413 / 256     & 319\,KiB $\le$ budget & --- \\
\bottomrule
\end{tabular}
\end{table}

\section{Reference Implementation: SciRust}
\label{sec:impl}

The kernel is written in Rust (2021 edition) with a no-allocation
invariant in the hot loops. Salient engineering properties:

\begin{itemize}[leftmargin=1.4em,itemsep=2pt]
  \item \textbf{Runtime SIMD dispatch.} The coarse dot product (128-dim
    INT4 dequant $+$ FMA) has scalar, AVX2, AVX-512 (x86\_64) and NEON
    (aarch64) paths, selected at run time via
    \texttt{is\_x86\_feature\_detected!} (order AVX-512 $>$ AVX2 $>$
    scalar) with a portable fallback. There is \emph{no} compile-time
    feature gating of correctness.
  \item \textbf{Branchless residual.} The Hamming term is a fixed
    four-word reduction with no data branches. It dispatches at run time to
    a single \texttt{vpopcntq} over all 256 residual bits
    (\texttt{\_mm256\_popcnt\_epi64}) on x86-64 CPUs advertising AVX-512
    \textsc{vpopcntdq}+\textsc{vl}, otherwise to \texttt{count\_ones()}
    ($\to$\,\texttt{POPCNT}/\texttt{CNT}). The two are asserted bit-equal on
    capable hardware and compile-checked elsewhere.
  \item \textbf{Zero dependencies.} The library has no runtime
    dependencies; \texttt{criterion} is a dev-only dependency for
    micro-benchmarks.
  \item \textbf{Tested equivalence.} 51 tests (unit, integration,
    property/fuzz, doctests, $\lambda$ calibration, Soft-Paging) assert,
    among other things, SIMD $\equiv$ scalar (randomized fuzz; AVX2,
    AVX-512, VPOPCNTDQ, NEON), the budget invariants of
    \texttt{enforce\_budget} over 300 randomized configurations
    ($\texttt{live\_bytes}\le\text{budget}$, consistent accounting, slot
    recycling), score finiteness, the Hamming identity, and conformance to
    Eq.~\eqref{eq:score}. The NEON and VPOPCNTDQ paths are checked by
    cross-compilation / feature-gated execution. Continuous integration runs
    \texttt{fmt}\,+\,\texttt{clippy -D warnings}\,+\,tests\,+\,bench
    compilation\,+\,AArch64 cross-compile.
\end{itemize}

\section{Experimental Evaluation}
\label{sec:eval}

\subsection{Setup and methodological honesty}
\label{sec:setup}
All numbers below are reproducible with fixed seeds. Measurements use
\emph{synthetic} data; the sign-LSH projection $Z$ is \emph{random} by
design. In \S\ref{sec:bincore}--\S\ref{sec:scorefp} the low-rank base is
assumed ideal ($W_{\text{down}}, W_{\text{up}}$ not learned) and we
isolate the INT4-quantization $+$ 1-bit-residual $+$ ranking machinery;
$\rho = \lVert e\rVert / \lVert k_{\text{real}}\rVert$ denotes the energy
fraction the base leaves to the residual.
\S\ref{sec:learned}--\S\ref{sec:taskaware} lift the ideal-base
assumption by learning the base (PCA, then task-aware SGD).
Throughput is reported on a \emph{shared} x86\_64 bench (treat as an
order of magnitude). We do \emph{not} report real-model perplexity or
hardware cache counters; \S\ref{sec:limits} explains why.

\subsection{Binary core fidelity (sign-LSH, $\ds=256$)}
\label{sec:bincore}
Over pairs spanning the full angular range, the 1-bit estimator tracks
the true cosine with Spearman $> 0.9$ (test-proven). In the
\emph{realistic} regime, where the residual is near-orthogonal to the
query, $256$ bits give only moderate resolution: Spearman(residual,
$\cos\theta) \approx 0.67$. This is a genuine limit---one bit per
hyperplane resolves nearly orthogonal directions poorly---and motivates
larger $\ds$ as future work.

\subsection{Full score vs.\ full-precision ground truth}
\label{sec:scorefp}
Table~\ref{tab:scorefp} reports rank fidelity of the full SLHA score
($N=512$ tokens, fixed query) against FP ground truth, as a function of
$\rho$.

\begin{table}[t]
\centering
\caption{Score rank fidelity vs.\ FP (ideal base). \textsc{hot}
$=$ latent $+$ residual; \textsc{warm} $=$ latent only.}
\label{tab:scorefp}
\small
\begin{tabular}{@{}cccccc@{}}
\toprule
$\rho$ & \textsc{hot} Spearman & \textsc{hot} top-16 & \textsc{warm} Spearman & \textsc{warm} top-16 \\
\midrule
0.05 & 0.987 & 0.875 & 0.987 & 0.875 \\
0.10 & 0.983 & 0.750 & 0.982 & 0.750 \\
0.20 & 0.966 & 0.750 & 0.961 & 0.688 \\
0.30 & 0.937 & 0.625 & 0.925 & 0.562 \\
0.50 & 0.844 & 0.500 & 0.811 & 0.438 \\
0.70 & 0.702 & 0.375 & 0.627 & 0.250 \\
\bottomrule
\end{tabular}
\end{table}

Two observations. (i)~At low $\rho \le 0.1$, \textsc{hot} $\approx$
\textsc{warm}: freeing the residual is near-lossless when the base
captures most of the energy---exactly the condition under which CCOS
moves a tile to \textsc{warm}. (ii)~The residual helps everywhere
(\textsc{hot} $\ge$ \textsc{warm}), with the gap widening as $\rho$ grows
($+0.08$ Spearman, $+12$ pts top-16 at $\rho=0.7$), but the gain is
modest at $256$ bits.

\subsection{Learned PCA base and grouped quantization}
\label{sec:learned}
Lifting the ideal-base assumption, the base is learned by PCA
(best rank-$\dc$ linear reconstruction, Eckart--Young~\cite{eckartyoung})
on synthetic keys with controllable spectrum; $Z$ stays random
(Table~\ref{tab:pca}).

\begin{table}[t]
\centering
\caption{Learned PCA base ($d_{\text{model}}=256$, latent $128$,
$\ds=256$); INT4 with grouped MX micro-scales.}
\label{tab:pca}
\small
\begin{tabular}{@{}cccccc@{}}
\toprule
decay & energy captured & \textsc{hot} Spearman & \textsc{hot} top-16 & \textsc{warm} Spearman & \textsc{warm} top-16 \\
\midrule
0.99 & 97.6\% & 0.788 & 0.562 & 0.703 & 0.438 \\
0.97 & 99.7\% & 0.814 & 0.375 & 0.700 & 0.312 \\
0.93 & 99.5\% & 0.884 & 0.438 & 0.609 & 0.188 \\
0.85 & 99.1\% & 0.905 & 0.562 & 0.871 & 0.500 \\
\bottomrule
\end{tabular}
\end{table}

\textsc{hot} plateaus at $0.79$--$0.90$ despite $97$--$99.7\%$ captured
energy: reconstruction energy is not score fidelity. \textsc{hot} $>$
\textsc{warm} everywhere, by as much as $+0.28$ Spearman (decay $0.93$).
Grouped MX cuts reconstruction error by $>2\times$ (test-proven) but the
end-to-end gain is marginal; whitening the latent \emph{degrades} it
($0.859 \to 0.750$). The plateau is examined in \S\ref{sec:ablation}.

\subsection{Attention-output fidelity (the headline result)}
\label{sec:output}
Rank fidelity is a proxy; what a model consumes is the attention
\emph{output} $\mathrm{out} = \mathrm{softmax}(QK^\top/\sqrt{d})\,V$.
Table~\ref{tab:output} reports cosine and relative-$L_2$ between the FP
output and the SLHA output (PCA base, $d=256$, $N=512$, mean over 64
queries, $d_v=64$).

\begin{table}[t]
\centering
\caption{Attention-output fidelity --- the closest accessible proxy for
perplexity drift.}
\label{tab:output}
\small
\begin{tabular}{@{}ccll@{}}
\toprule
decay & energy & \textsc{hot} $\cos$ / rel$L_2$ & \textsc{warm} $\cos$ / rel$L_2$ \\
\midrule
0.99 & 97.6\% & 0.948 / 0.318 & 0.943 / 0.333 \\
0.95 & 99.6\% & 0.989 / 0.147 & 0.988 / 0.154 \\
0.90 & 99.4\% & 0.994 / 0.108 & 0.993 / 0.114 \\
0.80 & 98.9\% & 0.997 / 0.078 & 0.996 / 0.082 \\
\bottomrule
\end{tabular}
\end{table}

This is the most important result. The output is far more robust than
the ranking suggested: where score Spearman plateaued at $0.79$--$0.90$,
the output cosine reaches $0.95$--$0.997$. The reason is that the softmax
averages values, so score errors between tokens of similar weight cancel.
\textsc{hot} $\ge$ \textsc{warm} throughout (small gap at these high
decays, where the residual matters little).

\subsection{Task-aware learned projection vs.\ PCA}
\label{sec:taskaware}
PCA picks the highest-variance directions of the \emph{keys} and ignores
the \emph{query} distribution. When $Q$ and $K$ favor different
subspaces, a projection trained to preserve the \emph{score}
$\inner{Q}{K}$ (not key reconstruction) wins. \texttt{train\_projection}
performs this descent with a closed-form gradient
($a=PQ$, $b=PK$, $r=\inner{Q}{K}-\inner{a}{b}$,
$\partial r^2/\partial P = -2r(bQ^\top + aK^\top)$), evaluated on an
adversarial case (orthonormal factors split into high key-variance /
high query-weight halves). We measure in \textsc{warm} to isolate the
projection (Table~\ref{tab:taskaware}).

\begin{table}[t]
\centering
\caption{Task-aware learned projection vs.\ PCA under deliberate $Q/K$
shift. SGD score loss $28.8 \to 5.1$.}
\label{tab:taskaware}
\small
\begin{tabular}{@{}lccc@{}}
\toprule
Projection & \textsc{warm} Spearman & \textsc{hot} Spearman & output $\cos$ \\
\midrule
PCA            & 0.161 & 0.325 & 0.947 \\
\textbf{Learned} & \textbf{0.859} & \textbf{0.863} & \textbf{0.985} \\
\bottomrule
\end{tabular}
\end{table}

Learning beats PCA decisively (\textsc{warm} $0.86$ vs.\ $0.16$): the
task-aware projection reallocates low-rank capacity toward the directions
that matter for the score. Caveats: synthetic data with deliberate $Q/K$
shift; when $Q$ and $K$ share statistics, PCA is already near-optimal;
and this is an SGD proof-of-concept on $P$ alone, not joint training with
a real model.

\subsection{Ablation: quantization is not the bottleneck}
\label{sec:ablation}
The latent can also be NF4 (normal-float codebook, same 4-bit budget,
same 128-byte tile), or TQ3---a port of the TurboQuant KV codec into the
same 64-byte latent: all 128 dims as 3-bit codes (48 bytes, symmetric
grid with no zero level) plus a separable per-dim 1-bit sign correction
(16 bytes, $\pm 0.25$ step). We add an INT8 reference (coarse only,
hypothetical 192-byte tile) to isolate bit-width
(Table~\ref{tab:ablation}).

\begin{table}[t]
\centering
\caption{Latent codec ablation at decay $0.93$.}
\label{tab:ablation}
\small
\begin{tabular}{@{}lcc@{}}
\toprule
Latent codec & \textsc{hot} Spearman & \textsc{warm} Spearman \\
\midrule
INT4 uniform (1 scale) & 0.879 & 0.603 \\
INT4 grouped (MX)      & 0.884 & 0.609 \\
NF4 grouped            & 0.885 & 0.610 \\
TQ3 (3-bit $+$ 1-bit corr.) & 0.881 & 0.598 \\
\textbf{INT8} (ref, $2\times$ bytes) & --- & \textbf{0.606} \\
\bottomrule
\end{tabular}
\end{table}

NF4 lowers reconstruction error (test-proven) but the end-to-end gain is
null-to-marginal ($0.884 \to 0.885$). TQ3 lands at the level of the
4-bit codecs (expected: $3+1$ bits of information per dim, same
$0.25$-step worst-case error) despite its zero-free grid costing
$1.3$--$1.6\times$ the MSE of grouped INT4 on a Gaussian latent
(test-bounded at $\le 2\times$); what it buys instead is separability:
dropping the 16-byte correction plane degrades gracefully to a pure
3-bit tile, a paging state finer than
\textsc{hot}$\to$\textsc{warm} (not yet wired into the cache manager).
Crucially, INT8 does \emph{not}
beat INT4 on \textsc{warm} ($\approx 0.61$): doubling the bits does not
lift the coarse ceiling. \textbf{The limiting factor is therefore the
low-rank projection, not latent quantization.} This corrects an earlier
hypothesis that blamed INT4, and refocuses the real levers on (a)~a
better (learned) projection and (b)~the 1-bit residual.

\subsection{$\lambda$ calibration}
\label{sec:calib}
\texttt{calibrate\_lambda} compares the residual weight to an FP
reference. Writing $\text{score} = \text{coarse} + \lambda r$ with
$r = \ds - 2\,\popc$ (independent of $\lambda$), the optimal multiplier
on the formula's $\lambda$ has the closed form
$\alpha^\star = \sum r_t (\lambda r) / \sum (\lambda r)^2$ with
$r_t = \inner{Q}{K} - \text{coarse}$ (Table~\ref{tab:calib}).

\begin{table}[t]
\centering
\caption{$\lambda$ calibration. $C_{\text{emp}} = \alpha^\star
\cdot \sqrt{\pi/(2\ds)}$.}
\label{tab:calib}
\small
\begin{tabular}{@{}cccccc@{}}
\toprule
$\rho$ & $\alpha^\star$ (LS) & $C_{\text{emp}}$ & RMSE @formula & RMSE @opt & $\Delta$out @formula \\
\midrule
0.10 & 4.18 & 0.327 & 1.37 & 1.26 & 0.006 \\
0.20 & 4.23 & 0.331 & 2.25 & 1.94 & 0.017 \\
0.30 & 4.24 & 0.332 & 3.29 & 2.79 & 0.036 \\
0.50 & 4.25 & 0.333 & 5.87 & 4.92 & 0.106 \\
0.70 & 4.26 & 0.334 & 9.90 & 8.26 & 0.262 \\
\bottomrule
\end{tabular}
\end{table}

The shape $\lambda \propto \sigE$ is validated ($\alpha^\star$ constant
at $4.18$--$4.26$). The analytic constant is $\approx 4.2\times$ too
small; the calibrated constant is $C_{\text{emp}} \approx 0.33$ at
$\ds=256$, which cuts score RMSE by $\sim 15\%$. We expose this value as
an opt-in (\texttt{LAMBDA\_C\_CALIBRATED}, with a test asserting its
optimal multiplier is $\approx 1$), but the crate keeps the analytic form
as default for two reasons: the $4.2\times$ factor is measured on
synthetic data and the true optimum is model-dependent; and calibration
minimizes score \emph{magnitude} (RMSE), a different objective from the
\emph{ranking} (top-$k$/Spearman) that attention ultimately consumes, so
over-weighting the directionally-noisy residual could trade ranking for
magnitude.

\subsection{Throughput (SIMD)}
\label{sec:throughput}
Each kernel path has an equivalence test against the scalar reference
(Table~\ref{tab:simd}).

\begin{table}[t]
\centering
\caption{Score throughput on a shared x86\_64 bench (order of
magnitude).}
\label{tab:simd}
\small
\begin{tabular}{@{}lcc@{}}
\toprule
Path & Throughput & Ratio \\
\midrule
Scalar (reference)          & $\sim$2.9 M scores/s  & $1\times$ \\
AVX2                        & $\sim$33.9 M scores/s & $\mathbf{11.5\times}$ \\
AVX-512 (16-wide FMA/group) & $\sim$41.6 M scores/s & $\mathbf{14.1\times}$ \\
\bottomrule
\end{tabular}
\end{table}

The speedup exceeds the nominal $8\times$/$16\times$ because the scalar
path also paid a branchy INT4 dequant that SIMD fuses. AVX-512 adds only
$\sim\!23\%$ over AVX2: the kernel is short (128 dims) and bound by
nibble-unpacking/loads, not FMA width. In reference cycles
(\texttt{rdtsc}): $\sim$942 cyc/tile scalar, $\sim$89 AVX2, $\sim$71
AVX-512. The NEON path is now \emph{measured} on a Jetson Thor AGX 128
(AArch64, Neoverse-V3AE) with the bench kit (Table~\ref{tab:thor}): NEON
sustains $17.1$ M scores/s vs.\ $3.0$ scalar, a $5.7\times$ device-level
speedup. SLHA v2 thus targets \emph{both} families---x86-64 servers and
AArch64 edge devices---and the crate ships a portable
\texttt{platform\_report} kit (and runner script) that detects features
(AVX/NEON/SVE2), reports the cache hierarchy, checks the tile size/alignment
against the device's line, and times throughput on whatever CPU runs it. The
x86-64 figures are the server baseline; the AArch64 figures are measured
on-device, not fabricated.

\begin{table}[t]
\centering
\caption{On-device measurement on a Jetson Thor AGX 128 (AArch64,
Neoverse-V3AE), via the \texttt{platform\_report} kit. All cache levels use
64-byte lines; \texttt{sve2} is present (cf.\ \S\ref{sec:future}).}
\label{tab:thor}
\small
\begin{tabular}{@{}ll@{}}
\toprule
Property & Value \\
\midrule
Cache lines (L1d / L1i / L2) & 64 B / 64 B / 64 B \\
SIMD features                & neon, dotprod, i8mm, fp16, sve, \textbf{sve2} \\
Tile size / align            & 128 B / 64 B (two 64-byte lines) \\
Dispatched kernel            & NEON \\
Throughput (NEON / scalar)   & $17.1$ / $3.0$ M scores/s ($\mathbf{5.7\times}$) \\
\bottomrule
\end{tabular}
\end{table}

\subsection{Memory traffic vs.\ a bf16 baseline}
\label{sec:membw}
\texttt{bench\_vs\_fp16} compares scoring an SLHA tile (128 B/token)
against a dot product on a bf16 key ($\dk \cdot 2 = 256$ B/token), both
in AVX2 (Table~\ref{tab:membw}).

\begin{table}[t]
\centering
\caption{SLHA (128 B/token) vs.\ bf16 (256 B/token), AVX2.}
\label{tab:membw}
\small
\begin{tabular}{@{}lccc@{}}
\toprule
Context (tokens) & footprint SLHA / bf16 & SLHA & bf16 \\
\midrule
8\,192      & 1 / 2\,MiB    & 42.4 M/s & 17.0 M/s \\
65\,536     & 8 / 16\,MiB   & 40.1 & 16.9 \\
262\,144    & 32 / 64\,MiB  & 35.4 & 14.0 \\
1\,048\,576 & 128 / 256\,MiB & 29.9 & 13.8 \\
\bottomrule
\end{tabular}
\end{table}

SLHA scores $\sim\!2.5\times$ more tokens/s at comparable memory
bandwidth: reading $2\times$ fewer bytes/token (plus a shorter INT4
unpack than bf16 decode) converts directly into throughput---the
bandwidth-wall thesis, verified at the kernel level. Throughput decays as
context grows (42 $\to$ 30 M/s), an indirect cache-residency effect.
\textbf{Honesty:} the LLC is 260\,MB on this bench, so we do not saturate
DRAM; the gain comes from byte volume and uops, not a DRAM limit. On a
smaller LLC or a genuinely DRAM-bound decode, SLHA's advantage would be
larger.

\section{Discussion: real levers and false levers}
\label{sec:discussion}
Synthesizing the evaluation:

\begin{itemize}[leftmargin=1.4em,itemsep=2pt]
  \item \textbf{Lever \#1 --- the low-rank projection.} A learned,
    task-aware projection beats PCA decisively under $Q/K$ shift
    (\textsc{warm} $0.16 \to 0.86$, \S\ref{sec:taskaware}).
  \item \textbf{Lever \#2 --- the 1-bit residual.} It recovers much of
    what the coarse term misses (\textsc{hot} $\gg$ \textsc{warm} at high
    $\rho$), once $\lambda$ is calibrated.
  \item \textbf{False lever --- latent bit-width.} INT8 does not beat
    INT4 on the coarse term; NF4 and MX grouping improve reconstruction
    but barely move end-to-end ranking (\S\ref{sec:ablation}).
  \item \textbf{False lever --- SIMD width.} AVX-512 adds only
    $\sim\!23\%$ over AVX2; the kernel is unpack/load-bound, not
    FMA-width-bound.
\end{itemize}

\noindent\textbf{Assumed negative results.} Latent whitening degrades
fidelity; MX/NF4 grouping yields marginal end-to-end gains; and our own
earlier conclusion that INT4 was the bottleneck was refuted by the INT8
measurement and corrected. We report these because they are as
informative as the positive results.

\section{Limitations and Threats to Validity}
\label{sec:limits}
We are explicit about the boundaries of these results.

\begin{itemize}[leftmargin=1.4em,itemsep=2pt]
  \item \textbf{Synthetic data, random $Z$.} All fidelity numbers use
    synthetic keys/queries and a random sign-LSH projection. They
    validate the \emph{mechanism}, not (yet) quality on a real model.
  \item \textbf{No real-model perplexity.} The closest proxy is the
    attention-output cosine (\S\ref{sec:output}); a real perplexity
    measurement requires a model and dataset not available in our
    sandbox.
  \item \textbf{No hardware cache counters.} \texttt{perf} is
    unavailable (\texttt{perf\_event\_paranoid=2}); cache effects are
    shown only indirectly via throughput decay with context size
    (\S\ref{sec:membw}).
  \item \textbf{Throughput on a shared bench.} SIMD ratios are an order
    of magnitude, not a controlled measurement; the LLC was large enough
    that DRAM was not saturated.
  \item \textbf{$\lambda$ constant is data-dependent.} The $4.2\times$
    correction is measured on synthetic data; the production default
    stays analytic for that reason.
  \item \textbf{Single ARM device.} ARM throughput is now measured, but on
    \emph{one} part (Jetson Thor AGX 128, Neoverse-V3AE; Table~\ref{tab:thor});
    other AArch64 microarchitectures may differ, and the same
    hardware-counter caveat applies.
  \item \textbf{Projection not jointly trained.} The task-aware result
    optimizes $P$ alone, not end-to-end with a transformer.
\end{itemize}

\section{Future Work}
\label{sec:future}
Two hardware directions we identified are now \emph{implemented} and described
above: the AVX-512 \textsc{vpopcntdq} residual reduction (\S\ref{sec:impl}),
and host-adaptive tile alignment via \texttt{build.rs} (item~2 below;
\S\ref{sec:cachemap}). The latter \emph{replaces} an earlier
architecture-conditional 128-byte alignment that we implemented and then
\emph{withdrew} once the device measurement (Table~\ref{tab:thor}) showed the
Jetson Thor uses 64-byte lines, not 128: instead of an \texttt{aarch64}-wide
assumption, the build script probes the actual host line size and upgrades only
where it helps. The remaining directions are proposals (not measured here).

\begin{enumerate}[leftmargin=1.6em,itemsep=2pt]
  \item \textbf{Scalable SIMD on ARM (Jetson Thor class).} The coarse dot is
    the hot term; an SVE2 path on Neoverse-V3AE---\texttt{sve2} is confirmed
    present on the measured device (Table~\ref{tab:thor})---could use
    \emph{scalable} vectors (adapting to the implementation's vector length)
    beyond fixed NEON. The residual popcount is already branchless and
    near-optimal at 256 bits, so the lever is the dot, not the popcount. We
    deliberately keep this on the roadmap rather than in the shipped kernel for
    a concrete toolchain reason, verified on the current stable compiler
    (\texttt{rustc 1.94.1}): runtime SVE2 \emph{detection}
    (\texttt{is\_aarch64\_feature\_detected!("sve2")}) is stable, but the SVE2
    \emph{intrinsics} (\texttt{svdot\_s32} and friends) are \emph{absent} from
    stable \texttt{core::arch::aarch64}---they, and portable \texttt{std::simd},
    remain nightly-gated---whereas the reference crate is intentionally stable
    and dependency-free. The one stable route, hand-written \texttt{asm!} with
    SVE2 mnemonics, \emph{compiles} but cannot be \emph{measured or
    test-checked} without an SVE2 device: our CI is x86-64, and cross-compiling
    to \texttt{aarch64} type-checks the Rust, not \texttt{asm!} semantics, while
    every other SIMD path in the crate is pinned by a bit-exact equivalence test
    (\S\ref{sec:tile})---so shipping an unverified kernel would break that
    contract. The NEON\,+\,\texttt{cnt} fallback remains correct meanwhile; SVE2
    will land behind runtime vector-length detection once the intrinsics
    stabilize on stable (or once an SVE2 device is available to validate an
    \texttt{asm!} kernel against the scalar reference) --- and requires ARM
    hardware to measure.
  \item \textbf{Host-adaptive alignment via \texttt{build.rs} (implemented).}
    The tile defaults to \texttt{align(64)}, correct for every 64-byte-line
    part (all our targets, including the Thor). The only case that benefits
    from \texttt{align(128)} is a part with genuine 128-byte lines (e.g.\ Apple
    Silicon). Since alignment is a compile-time constant, the shipped
    \texttt{build.rs} probes the host's \emph{actual} L1d line size on a
    \emph{native} build (host triple $=$ target triple; Linux \texttt{sysfs} or
    macOS \texttt{sysctl}) and emits \texttt{cfg(cache\_line\_128)} to raise the
    alignment only where it helps---never the \emph{wrong} ``AArch64
    $\Rightarrow$ 128'' assumption we measured and discarded. For
    cross-compilation the host line size is irrelevant, so the safe 64-byte
    default holds. A portability refinement, not a correctness change: the tile
    stays 128~bytes with zero padding either way, and on every one of our
    targets the result is unchanged (\texttt{align(64)}).
  \item \textbf{Salient-outlier tiles (BiLLM-style).} Following
    BiLLM~\cite{billm}, isolating a few \emph{salient} dimensions in high
    precision (a \texttt{salient\_mask} plus a small FP block) while
    binarizing the rest could raise fidelity within the 128-byte budget. Two
    non-obvious issues make this a design study rather than a drop-in change.
    \emph{(i)~Metadata budget:} a 16-byte salient block must displace
    existing metadata, and the natural victims---$\sigE$ and the MX
    \texttt{group\_scales}---are exactly the fields that drive CCOS paging
    (\S\ref{sec:ccos}) and microscaling, so the trade-off must be measured,
    not assumed. \emph{(ii)~Expected benefit --- now measured.} Our ablation
    (\S\ref{sec:ablation}) shows quantization is not the bottleneck on Gaussian
    keys; BiLLM targets a different regime, so we extended the benchmark with
    injected channel outliers (the \texttt{salient\_outliers} example). The
    mechanism is real at the reconstruction level: with outlier channels of
    $16$--$32\times$ magnitude, INT4 latent RMSE blows up ($0.08\!\to\!0.52$)
    while salient-$s$ stays flat \emph{provided} $s$ covers the outlier channels
    ($s{=}4$ keeps RMSE $\approx 0.08$ throughout). But two caveats temper it:
    the attention \emph{output} under plain INT4 stays at cosine $\ge 0.977$
    even at $32\times$ (the softmax absorbs the error, cf.\
    \S\ref{sec:output}), so the end-task gain is modest; and the tile's
    affordable budget ($s{=}2$ FP values) can \emph{under}-perform plain INT4
    when outliers outnumber it (partial preservation imbalances the remaining
    scale). Measured verdict: worthwhile only if a target model's
    outlier-channel count is known to fit the budget---otherwise the 16 bytes
    are better spent on $\sigE$ and the MX scales. A \emph{measured} defer, not
    a guess.
\end{enumerate}

Beyond these: larger $\ds$ to lift the near-orthogonal resolution limit
(\S\ref{sec:bincore}); \textbf{joint} training of $W_{\text{down}},
W_{\text{up}}$ with a real model (\S\ref{sec:taskaware} establishes the
principle); and real hardware validation (\texttt{perf} cache counters,
end-to-end decode throughput, perplexity).

\section{Conclusion}
\label{sec:conclusion}
SLHA~v2 marries a bit-exact, cache-aware data structure with a hybrid
low-rank/1-bit attention score so that the KV cache becomes a fluid,
architecture-aware structure rather than a monolithic burden. The
mechanism is mathematically correct (test-proven, including
scalar/AVX2/AVX-512 equivalence) and directionally validated:
\textsc{hot} $\ge$ \textsc{warm} everywhere; Soft-Paging is near-lossless
at low residual energy (output cosine $\approx 0.9995$ when half the
working set is paged); SIMD reaches $\sim\!13\times$;
SLHA scores $\sim\!2.5\times$ more tokens/s than bf16 at the kernel
level; a learned projection beats PCA under $Q/K$ shift; and the
attention output stays at cosine $0.95$--$0.997$ versus full precision. A
controlled ablation reframes the problem: the coarse-score ceiling is set
by the low-rank projection, not by latent quantization. What remains is
real-hardware and real-model validation. We release the reference
implementation, benchmarks, and tests to make every claim reproducible.

\paragraph{Reproducibility.}
The reference crate (\textsc{SciRust}), all benchmarks, and the 51 tests
backing the numbers in \S\ref{sec:eval} are available in the project
repository; every measurement uses fixed random seeds.

\begin{thebibliography}{99}
\setlength{\itemsep}{2pt}

\bibitem{vaswani} A.~Vaswani, N.~Shazeer, N.~Parmar, J.~Uszkoreit,
L.~Jones, A.~N.~Gomez, {\L}.~Kaiser, and I.~Polosukhin.
\newblock Attention Is All You Need.
\newblock In \emph{NeurIPS}, 2017.

\bibitem{deepseekv2} DeepSeek-AI.
\newblock DeepSeek-V2: A Strong, Economical, and Efficient
Mixture-of-Experts Language Model.
\newblock arXiv:2405.04434, 2024.

\bibitem{vllm} W.~Kwon, Z.~Li, S.~Zhuang, Y.~Sheng, L.~Zheng,
C.~H.~Yu, J.~E.~Gonzalez, H.~Zhang, and I.~Stoica.
\newblock Efficient Memory Management for Large Language Model Serving
with PagedAttention.
\newblock In \emph{SOSP}, 2023.

\bibitem{charikar} M.~S.~Charikar.
\newblock Similarity Estimation Techniques from Rounding Algorithms.
\newblock In \emph{STOC}, 2002.

\bibitem{jl} W.~B.~Johnson and J.~Lindenstrauss.
\newblock Extensions of Lipschitz Mappings into a Hilbert Space.
\newblock \emph{Contemporary Mathematics}, 26:189--206, 1984.

\bibitem{eckartyoung} C.~Eckart and G.~Young.
\newblock The Approximation of One Matrix by Another of Lower Rank.
\newblock \emph{Psychometrika}, 1(3):211--218, 1936.

\bibitem{qlora} T.~Dettmers, A.~Pagnoni, A.~Holtzman, and L.~Zettlemoyer.
\newblock QLoRA: Efficient Finetuning of Quantized LLMs.
\newblock In \emph{NeurIPS}, 2023.

\bibitem{gptq} E.~Frantar, S.~Ashkboos, T.~Hoefler, and D.~Alistarh.
\newblock GPTQ: Accurate Post-Training Quantization for Generative
Pre-trained Transformers.
\newblock In \emph{ICLR}, 2023.

\bibitem{awq} J.~Lin, J.~Tang, H.~Tang, S.~Yang, X.~Dang, and S.~Han.
\newblock AWQ: Activation-aware Weight Quantization for LLM Compression
and Acceleration.
\newblock In \emph{MLSys}, 2024.

\bibitem{billm} W.~Huang, Y.~Liu, H.~Qin, Y.~Li, S.~Zhang, X.~Liu,
M.~Magno, and X.~Qi.
\newblock BiLLM: Pushing the Limit of Post-Training Quantization for
LLMs.
\newblock In \emph{ICML}, 2024.

\bibitem{kivi} Z.~Liu, J.~Yuan, H.~Jin, S.~Zhong, Z.~Xu, V.~Braverman,
B.~Chen, and X.~Hu.
\newblock KIVI: A Tuning-Free Asymmetric 2bit Quantization for KV Cache.
\newblock In \emph{ICML}, 2024.

\bibitem{h2o} Z.~Zhang, Y.~Sheng, T.~Zhou, T.~Chen, L.~Zheng, R.~Cai,
Z.~Song, Y.~Tian, C.~R\'e, C.~Barrett, Z.~Wang, and B.~Chen.
\newblock H2O: Heavy-Hitter Oracle for Efficient Generative Inference of
Large Language Models.
\newblock In \emph{NeurIPS}, 2023.

\bibitem{streamingllm} G.~Xiao, Y.~Tian, B.~Chen, S.~Han, and M.~Lewis.
\newblock Efficient Streaming Language Models with Attention Sinks.
\newblock In \emph{ICLR}, 2024.

\end{thebibliography}

\end{document}
